% =====================================================================
% chapters/unitary.tex　第十二章　酉空间
% 由 document.tex 在第十一章（矩阵分解）之后、第二部分综合项目之前 \input.
% =====================================================================
	\chapter{酉空间}\label{ch:unitary}
	
	\paragraph*{本章地位}
	
	第十章的内积只在实数域上定义. 但复数在第八章已经不可避免：实矩阵的特征值可以是复数，旋转矩阵就没有实特征向量. 本章把内积推广到复数域，得到酉空间，并把第十章的正交矩阵、第十一章的谱定理与奇异值分解推广到复矩阵. 核心结论有两个：每个复方阵都能用酉矩阵化为上三角矩阵（Schur 分解）；一个复方阵能用酉矩阵对角化，当且仅当它是正规矩阵. 推广之后，实数情形的许多``例外''消失了，结论反而更加整齐.
	
	\chapterhours{—}{—}{6}
	
	\begin{introduction}[学习目标]
		\item 说明复内积为什么要取共轭，验证给定的函数是否为复内积
		\item 在酉空间中作 Gram–Schmidt 正交化，求向量在标准正交基下的坐标
		\item 求复矩阵的共轭转置与线性变换的伴随变换，判断矩阵是否为酉矩阵、Hermite 矩阵、正规矩阵
		\item 用 Schur 分解与正规矩阵的谱定理把复矩阵酉上三角化或酉对角化
		\item 说明实对称矩阵、正交矩阵的结论是本章结论的特例
	\end{introduction}
	
	\paragraph*{前置知识}
	
	第四章（线性空间与线性变换）；第八章（特征值、相似对角化）；第九章（Jordan 标准形，只在个别注记中用到）；第十章（内积、Gram–Schmidt 正交化、正交矩阵）；第十一章（谱定理、奇异值分解）；附录~A 中的复数.
	
	\paragraph*{学习提示}
	\begin{itemize}
		\item 本章的大多数定义与定理都能在第十、十一章找到实数版本. 学习时应把两者并排比较：哪些地方只需把转置换成共轭转置，哪些地方必须重新论证.
		\item 取共轭是本章唯一的新运算，也是最容易出错的地方. 每遇到一个公式，都应弄清共轭落在哪个量上.
		\item 本章的两个分解定理是全章的终点. Schur 分解对一切复方阵成立，可以作为第九章 Jordan 标准形在数值计算中的替代.
	\end{itemize}
	
	\newpage
	
	\begin{tcolorbox}[epigraph]
		\begin{quote}It has been written that the shortest and best way between two truths of the real domain often passes through the imaginary one.\end{quote}
		\begin{flushright}——Jacques Hadamard\end{flushright}
		\begin{quote}有人写道：实数领域中两个真理之间最短、最好的道路，往往要经过虚数领域.\end{quote}
		\begin{flushright}——雅克·阿达马\footnotemark\end{flushright}
	\end{tcolorbox}
	\footnotetext{雅克·阿达马（1865—1963），法国数学家，1896 年证明了素数定理. 引文出自他的《数学领域中的发明心理学》（1945）. 本章的正规矩阵就是一个例子：一个实矩阵能否用``旋转''把它化简，答案要到复数域中去找.}
	
	\section{\hC 复内积空间}\label{sec:complexip}
	
	先看照搬第十章的定义会出什么问题. 在 $\mathbb{C}^2$ 中取 $\vv x=(1,\mathrm i)^\trans$，若仍令 $\langle\vv x,\vv x\rangle=\vv x^\trans\vv x$，则
	\[\vv x^\trans\vv x=1^2+\mathrm i^2=0.\]
	一个非零向量的``长度''为零，长度就失去了意义. 补救的办法是取共轭：对复数 $z$，$z\bar z=|z|^2\ge0$，所以
	\[x_1\bar x_1+x_2\bar x_2+\cdots+x_n\bar x_n=|x_1|^2+|x_2|^2+\cdots+|x_n|^2\]
	是非负实数，且只在 $\vv x=\zero$ 时为零. 代价是对称性要改为``共轭对称''.
	
	\begin{definition}\label{def:complexip}
		设 $V$ 是复数域上的线性空间. 若对任意 $\vv\alpha,\vv\beta\in V$ 有一个复数 $\langle\vv\alpha,\vv\beta\rangle$ 与之对应，且满足：对一切 $\vv\alpha,\vv\beta,\vv\gamma\in V$，$k\in\mathbb{C}$，
		\begin{enumerate}
			\item （共轭对称性）$\langle\vv\alpha,\vv\beta\rangle=\overline{\langle\vv\beta,\vv\alpha\rangle}$；
			\item （对第一个变量线性）$\langle\vv\alpha+\vv\beta,\vv\gamma\rangle=\langle\vv\alpha,\vv\gamma\rangle+\langle\vv\beta,\vv\gamma\rangle$，$\langle k\vv\alpha,\vv\beta\rangle=k\langle\vv\alpha,\vv\beta\rangle$；
			\item （正定性）$\langle\vv\alpha,\vv\alpha\rangle\ge0$，且等号成立当且仅当 $\vv\alpha=\zero$，
		\end{enumerate}
		则称 $\langle\cdot,\cdot\rangle$ 是 $V$ 上的一个\textbf{复内积}\index{fu4 nei4ji1@复内积}，配备了复内积的复线性空间称为\textbf{复内积空间}；有限维的复内积空间称为\textbf{酉空间}\index{you3 kong1jian1@酉空间}(unitary space).
	\end{definition}
	
	由共轭对称性，$\langle\vv\alpha,\vv\alpha\rangle=\overline{\langle\vv\alpha,\vv\alpha\rangle}$ 总是实数，所以正定性中的``$\ge0$''是有意义的. 对第二个变量，内积是\textbf{共轭线性}的：
	\[\langle\vv\alpha,k\vv\beta\rangle=\overline{\langle k\vv\beta,\vv\alpha\rangle}=\overline{k\langle\vv\beta,\vv\alpha\rangle}=\bar k\langle\vv\alpha,\vv\beta\rangle.\]
	
	\begin{pitfallbox}
		$\langle\vv\alpha,k\vv\beta\rangle=\bar k\langle\vv\alpha,\vv\beta\rangle$，\textbf{不是} $k\langle\vv\alpha,\vv\beta\rangle$. 数从第二个位置提出来时要取共轭. 另外，物理学中常用相反的约定：内积对第二个变量线性、对第一个变量共轭线性（记为 $\langle\alpha|\beta\rangle$）. 阅读其他书籍时应先弄清作者用的是哪一种.
	\end{pitfallbox}
	
	\begin{example}
		\begin{enumerate}
			\item $\mathbb{C}^n$ 上的\textbf{标准内积}
			\[\langle\vv x,\vv y\rangle=x_1\bar y_1+x_2\bar y_2+\cdots+x_n\bar y_n.\]
			本章提到 $\mathbb{C}^n$ 时，若不加说明，都配备这个内积. 当 $\vv x,\vv y$ 都是实向量时，它就是第十章的标准内积.
			\item 闭区间 $[a,b]$ 上的全体复值连续函数（实部、虚部都连续）构成复线性空间，$\langle f,g\rangle=\int_a^bf(x)\overline{g(x)}\,\mathrm dx$ 是它上面的复内积.
			\item 全体 $m\times n$ 复矩阵构成的线性空间上，$\langle\vv A,\vv B\rangle=\sum_{i,j}a_{ij}\bar b_{ij}$，即把矩阵``拉直''后的标准内积.
		\end{enumerate}
	\end{example}
	
	为了用矩阵书写标准内积，引入共轭转置.
	
	\begin{definition}
		把复矩阵 $\vv A=(a_{ij})$ 的每个元素取共轭，得到的矩阵记为 $\bar{\vv A}=(\bar a_{ij})$. 矩阵 $\vv A^\herm=\bar{\vv A}^\trans$ 称为 $\vv A$ 的\textbf{共轭转置}\index{gong4e4 zhuan3zhi4@共轭转置}(conjugate transpose).
	\end{definition}
	
	于是 $\mathbb{C}^n$ 的标准内积可以写成
	\[\langle\vv x,\vv y\rangle=\vv y^\herm\vv x.\]
	注意 $\vv y$ 写在左边：这样内积才对第一个变量 $\vv x$ 线性. 共轭转置的运算规则与转置相同，只是数乘要取共轭：
	\[(\vv A+\vv B)^\herm=\vv A^\herm+\vv B^\herm,\quad(k\vv A)^\herm=\bar k\vv A^\herm,\quad(\vv{AB})^\herm=\vv B^\herm\vv A^\herm,\quad(\vv A^\herm)^\herm=\vv A.\]
	对实矩阵，$\vv A^\herm=\vv A^\trans$. 本书已经用 $\vv A^*$ 表示伴随矩阵（第五章），所以共轭转置记为 $\vv A^\herm$，不记为 $\vv A^*$.
	
	长度、单位向量与正交的定义与第十章相同：$\|\vv\alpha\|=\sqrt{\langle\vv\alpha,\vv\alpha\rangle}$；$\langle\vv\alpha,\vv\beta\rangle=0$ 时称 $\vv\alpha$ 与 $\vv\beta$ 正交. 由共轭对称性，$\langle\vv\alpha,\vv\beta\rangle=0$ 与 $\langle\vv\beta,\vv\alpha\rangle=0$ 等价，正交关系仍是对称的. 复内积空间中没有``夹角''：$\langle\vv\alpha,\vv\beta\rangle$ 是复数，不能等于一个余弦.
	
	\begin{theorem}[Cauchy–Schwarz 不等式]\label{thm:cscomplex}
		在复内积空间中，对任意 $\vv\alpha,\vv\beta$，
		\[|\langle\vv\alpha,\vv\beta\rangle|\le\|\vv\alpha\|\,\|\vv\beta\|,\]
		等号成立当且仅当 $\vv\alpha,\vv\beta$ 线性相关.
	\end{theorem}
	\begin{proof}
		$\vv\beta=\zero$ 时显然. 设 $\vv\beta\ne\zero$. 第十章的证明用到``实系数二次函数的判别式''，这里 $t$ 要取复数，不能照搬. 改为直接取 $t=\dfrac{\langle\vv\alpha,\vv\beta\rangle}{\|\vv\beta\|^2}$，展开
		\[0\le\langle\vv\alpha-t\vv\beta,\vv\alpha-t\vv\beta\rangle=\|\vv\alpha\|^2-\bar t\langle\vv\alpha,\vv\beta\rangle-t\langle\vv\beta,\vv\alpha\rangle+|t|^2\|\vv\beta\|^2.\]
		后三项都等于 $\dfrac{|\langle\vv\alpha,\vv\beta\rangle|^2}{\|\vv\beta\|^2}$（注意 $\langle\vv\beta,\vv\alpha\rangle=\overline{\langle\vv\alpha,\vv\beta\rangle}$），所以右边等于 $\|\vv\alpha\|^2-\dfrac{|\langle\vv\alpha,\vv\beta\rangle|^2}{\|\vv\beta\|^2}$，即得不等式. 等号成立当且仅当 $\vv\alpha=t\vv\beta$.
	\end{proof}
	
	这里的 $t\vv\beta$ 正是 $\vv\alpha$ 在 $\vv\beta$ 方向上的投影. 由 Cauchy–Schwarz 不等式，与第十章完全相同地得到三角不等式 $\|\vv\alpha+\vv\beta\|\le\|\vv\alpha\|+\|\vv\beta\|$ 与勾股定理.
	
	在实内积空间中，内积可以用长度表示（定理~\ref{thm:orthequiv} 的证明）. 在复内积空间中，展开 $\|\vv\alpha+c\vv\beta\|^2=\|\vv\alpha\|^2+\bar c\langle\vv\alpha,\vv\beta\rangle+c\langle\vv\beta,\vv\alpha\rangle+|c|^2\|\vv\beta\|^2$，分别取 $c=1,\mathrm i,-1,-\mathrm i$ 并组合，可以得到\textbf{极化恒等式}\index{ji2hua4 heng2deng3shi4@极化恒等式}
	\begin{equation}\label{eq:polarization}
		\langle\vv\alpha,\vv\beta\rangle=\frac14\sum_{k=0}^{3}\mathrm i^k\,\|\vv\alpha+\mathrm i^k\vv\beta\|^2.
	\end{equation}
	实数情形只需要两项，复数情形需要四项：两项确定实部，两项确定虚部.
	
	\section{\hC 正交性与标准正交基}\label{sec:unitaryonb}
	
	正交向量组、标准正交向量组、标准正交基的定义与第十章相同. 第十章关于它们的论证，只用到内积的线性与正定性，在酉空间中都成立，只是要注意共轭落在哪里.
	
	\begin{proposition}\label{prop:unitarycoord}
		设 $\vv e_1,\ldots,\vv e_n$ 是酉空间 $V$ 的标准正交基. 则对任意 $\vv\alpha,\vv\beta\in V$，
		\[\vv\alpha=\sum_{i=1}^n\langle\vv\alpha,\vv e_i\rangle\vv e_i,\qquad \langle\vv\alpha,\vv\beta\rangle=\sum_{i=1}^n\langle\vv\alpha,\vv e_i\rangle\overline{\langle\vv\beta,\vv e_i\rangle},\qquad\|\vv\alpha\|^2=\sum_{i=1}^n|\langle\vv\alpha,\vv e_i\rangle|^2.\]
		换言之，在标准正交基下，向量的内积等于坐标的标准内积.
	\end{proposition}
	\begin{proof}
		设 $\vv\alpha=\sum x_i\vv e_i$，与 $\vv e_j$ 作内积，由对第一个变量的线性与正交性，$\langle\vv\alpha,\vv e_j\rangle=x_j$. 设 $\vv\beta=\sum y_i\vv e_i$，则 $\langle\vv\alpha,\vv\beta\rangle=\sum_{i,j}x_i\bar y_j\langle\vv e_i,\vv e_j\rangle=\sum_ix_i\bar y_i$. 第三式是第二式的特例.
	\end{proof}
	
	第三式称为 \textbf{Parseval 等式}\index{parseval deng3shi4@Parseval 等式}. 若 $\vv e_1,\ldots,\vv e_m$ 只是标准正交向量组（不一定是基），同样的计算给出 \textbf{Bessel 不等式} $\sum_{i=1}^m|\langle\vv\alpha,\vv e_i\rangle|^2\le\|\vv\alpha\|^2$：把 $\vv\alpha$ 写成 $\sum_{i\le m}\langle\vv\alpha,\vv e_i\rangle\vv e_i$ 与一个和各 $\vv e_i$ 都正交的向量之和，再用勾股定理.
	
	\paragraph*{Gram–Schmidt 正交化.} 定理~\ref{thm:gramschmidt} 的公式与证明在酉空间中原样成立：
	\[\vv\beta_k=\vv\alpha_k-\sum_{i=1}^{k-1}\frac{\langle\vv\alpha_k,\vv\beta_i\rangle}{\langle\vv\beta_i,\vv\beta_i\rangle}\vv\beta_i.\]
	验证 $\langle\vv\beta_k,\vv\beta_j\rangle=0$ 时，系数从第一个位置提出，不取共轭. 所以系数必须是 $\langle\vv\alpha_k,\vv\beta_i\rangle$，而不是 $\langle\vv\beta_i,\vv\alpha_k\rangle$——在实数情形二者相等，在复数情形相差一个共轭. 由此，酉空间都有标准正交基，任一标准正交向量组都能扩充为标准正交基.
	
	\begin{example}
		在 $\mathbb{C}^3$ 中把 $\vv\alpha_1=(1,\mathrm i,0)^\trans$，$\vv\alpha_2=(\mathrm i,0,1)^\trans$ 正交化.
	\end{example}
	\begin{solution}
		$\vv\beta_1=\vv\alpha_1$，$\langle\vv\beta_1,\vv\beta_1\rangle=1+|\mathrm i|^2=2$. $\langle\vv\alpha_2,\vv\beta_1\rangle=\mathrm i\cdot\bar1+0\cdot\bar{\mathrm i}+1\cdot\bar0=\mathrm i$，所以
		\[\vv\beta_2=\vv\alpha_2-\frac{\mathrm i}2\vv\beta_1=\Big(\frac{\mathrm i}2,\ \frac12,\ 1\Big)^\trans.\]
		验证：$\langle\vv\beta_2,\vv\beta_1\rangle=\frac{\mathrm i}2\cdot1+\frac12\cdot(-\mathrm i)+0=0$. 单位化，$\|\vv\beta_1\|=\sqrt2$，$\|\vv\beta_2\|=\sqrt{\frac14+\frac14+1}=\sqrt{\frac32}$，得
		\[\vv e_1=\frac1{\sqrt2}(1,\mathrm i,0)^\trans,\qquad\vv e_2=\frac1{\sqrt6}(\mathrm i,1,2)^\trans.\]
		若误把系数取成 $\langle\vv\beta_1,\vv\alpha_2\rangle=-\mathrm i$，得到的 $\vv\alpha_2+\frac{\mathrm i}2\vv\beta_1=(\frac{3\mathrm i}2,-\frac12,1)^\trans$ 与 $\vv\beta_1$ 的内积是 $2\mathrm i\ne0$.
	\end{solution}
	
	\paragraph*{正交补与正交投影.} 子空间 $W$ 的正交补 $W^\perp$ 的定义与第十章相同. 取 $W$ 的标准正交基 $\vv e_1,\ldots,\vv e_k$ 并扩充为 $V$ 的标准正交基，与定理~\ref{thm:orthdecomp} 完全相同地得到
	\[V=W\oplus W^\perp,\qquad(W^\perp)^\perp=W,\]
	且 $\vv\alpha$ 在 $W$ 上的正交投影为 $\sum_{i=1}^k\langle\vv\alpha,\vv e_i\rangle\vv e_i$，它是 $W$ 中与 $\vv\alpha$ 距离最近的向量.
	
	\subsection{离散 Fourier 基}
	
	下面是酉空间中最重要的一组标准正交基. 记 $\omega=\mathrm e^{2\pi\mathrm i/n}$，它满足 $\omega^n=1$，$\bar\omega=\omega^{-1}$. 在 $\mathbb{C}^n$ 中令
	\[\vv f_k=\frac1{\sqrt n}\big(1,\ \omega^k,\ \omega^{2k},\ \ldots,\ \omega^{(n-1)k}\big)^\trans\qquad(k=0,1,\ldots,n-1).\]
	
	\begin{proposition}\label{prop:fourierbasis}
		$\vv f_0,\vv f_1,\ldots,\vv f_{n-1}$ 是 $\mathbb{C}^n$ 的标准正交基.
	\end{proposition}
	\begin{proof}
		$\langle\vv f_j,\vv f_k\rangle=\frac1n\sum_{m=0}^{n-1}\omega^{jm}\overline{\omega^{km}}=\frac1n\sum_{m=0}^{n-1}q^m$，其中 $q=\omega^{j-k}$. 若 $j=k$，$q=1$，和为 $n$，内积为 $1$. 若 $j\ne k$，由 $0<|j-k|<n$ 知 $q\ne1$，而 $q^n=1$，由等比数列求和公式，$\sum_{m=0}^{n-1}q^m=\frac{q^n-1}{q-1}=0$. 所以 $n$ 个向量两两正交且都是单位向量，构成标准正交基.
	\end{proof}
	
	向量 $\vv x\in\mathbb{C}^n$ 在这组基下的坐标是
	\[c_k=\langle\vv x,\vv f_k\rangle=\frac1{\sqrt n}\sum_{m=0}^{n-1}x_m\,\mathrm e^{-2\pi\mathrm ikm/n}\qquad(k=0,\ldots,n-1),\]
	这里把 $\vv x$ 的分量编号为 $x_0,\ldots,x_{n-1}$. $(c_0,\ldots,c_{n-1})$ 称为 $\vv x$ 的\textbf{离散 Fourier 变换}\index{li2san4 fourier bian4huan4@离散 Fourier 变换}. 由命题~\ref{prop:unitarycoord}，$\vv x=\sum c_k\vv f_k$，且 $\sum|x_m|^2=\sum|c_k|^2$. 直观地说，$\vv f_k$ 是``每 $n$ 个时刻转 $k$ 圈''的振动，坐标 $c_k$ 衡量 $\vv x$ 中这种振动的成分有多大.
	
	\begin{linkbox}{季节性与周期}
		月度的零售额、用电量、客运量等经济数据，常常每 $12$ 个月重复一次相似的起伏，称为\textbf{季节性}. 把 $n$ 个月的数据看作 $\mathbb{C}^n$ 中的向量，在 Fourier 基下展开，季节性就集中在``每 $12$ 个月转一圈''的那几个频率上. 统计部门发布的``季节调整后''数据，基本想法就是从数据中去掉这些分量，让``比上月增长多少''这样的比较不再受春节、暑期等因素干扰. 本章末的建模实验会把这一想法做成一个简单的模型.
	\end{linkbox}
	
	\begin{codebox}{Fourier 矩阵}
		\mstep{目标} 用矩阵验证 Fourier 基的标准正交性，并与快速 Fourier 变换的库函数对照.
		
		\mstep{准备} \texttt{numpy}；\texttt{numpy.fft.fft}.
		
		\mstep{任务}
		\begin{enumerate}
			\item 取 $n=8$，构造以 $\vv f_0,\ldots,\vv f_7$ 为列的矩阵 $\vv F$，计算 $\vv F^\herm\vv F$（\texttt{F.conj().T @ F}），与单位矩阵比较.
			\item 随机取 $\vv x\in\mathbb{C}^8$，计算坐标 $\vv F^\herm\vv x$，与 \texttt{numpy.fft.fft(x)} 比较，找出二者相差的常数因子.
			\item 验证 Parseval 等式：$\|\vv x\|^2=\|\vv F^\herm\vv x\|^2$.
		\end{enumerate}
		
		\mstep{检查点} $\vv F^\herm\vv F=\vv I$（误差在 $10^{-15}$ 量级）；\texttt{fft(x)} 等于 $\sqrt8\,\vv F^\herm\vv x$——库函数省略了归一化系数.
		
		\mstep{思考} 直接计算 $\vv F^\herm\vv x$ 需要约 $n^2$ 次乘法，\texttt{fft} 只需约 $n\log_2n$ 次. 当 $n=2^{20}$ 时二者相差多少倍？
	\end{codebox}
	
	\section{\hC 伴随变换与共轭转置}\label{sec:adjoint}
	
	第十一章谱定理的证明中，对称矩阵的关键性质是 $\langle\vv{Ax},\vv y\rangle=\langle\vv x,\vv{Ay}\rangle$. 在 $\mathbb{C}^n$ 中，对任意复方阵 $\vv A$，
	\[\langle\vv{Ax},\vv y\rangle=\vv y^\herm\vv{Ax}=(\vv A^\herm\vv y)^\herm\vv x=\langle\vv x,\vv A^\herm\vv y\rangle.\]
	即 $\vv A$ 从内积的第一个位置``挪到''第二个位置，就变成 $\vv A^\herm$. 把这件事推广到一般的线性变换：
	
	\begin{definition}
		设 $\mathscr T$ 是酉空间 $V$ 上的线性变换. 若 $V$ 上的线性变换 $\mathscr T^*$ 满足
		\[\langle\mathscr T\vv\alpha,\vv\beta\rangle=\langle\vv\alpha,\mathscr T^*\vv\beta\rangle\qquad\text{对一切}\ \vv\alpha,\vv\beta\in V,\]
		称 $\mathscr T^*$ 为 $\mathscr T$ 的\textbf{伴随变换}\index{ban4sui2 bian4huan4@伴随变换}(adjoint). 若 $\mathscr T^*=\mathscr T$，称 $\mathscr T$ 为\textbf{自伴变换}.
	\end{definition}
	
	要证明伴随变换存在，先看酉空间上的线性函数（第四章：$V\to\mathbb{C}$ 的线性映射）长什么样.
	
	\begin{lemma}\label{lem:riesz}
		设 $f$ 是酉空间 $V$ 上的线性函数，则存在唯一的 $\vv\gamma\in V$，使 $f(\vv\alpha)=\langle\vv\alpha,\vv\gamma\rangle$ 对一切 $\vv\alpha$ 成立.
	\end{lemma}
	\begin{proof}
		取标准正交基 $\vv e_1,\ldots,\vv e_n$，令 $\vv\gamma=\sum_i\overline{f(\vv e_i)}\,\vv e_i$. 对 $\vv\alpha=\sum x_i\vv e_i$，由命题~\ref{prop:unitarycoord}，$\langle\vv\alpha,\vv\gamma\rangle=\sum_ix_i\overline{\overline{f(\vv e_i)}}=\sum_ix_if(\vv e_i)=f(\vv\alpha)$. 唯一性：若 $\langle\vv\alpha,\vv\gamma\rangle=\langle\vv\alpha,\vv\gamma'\rangle$ 对一切 $\vv\alpha$ 成立，取 $\vv\alpha=\vv\gamma-\vv\gamma'$，得 $\|\vv\gamma-\vv\gamma'\|^2=0$.
	\end{proof}
	
	这个引理说，在酉空间中，``线性函数''与``向量''可以一一对应起来. 第十四章会看到，没有内积时，线性函数构成另一个空间（对偶空间），不能与 $V$ 自然地等同.
	
	\begin{theorem}\label{thm:adjoint}
		酉空间 $V$ 上的每个线性变换 $\mathscr T$ 都有唯一的伴随变换 $\mathscr T^*$. 若 $\mathscr T$ 在标准正交基 $\vv e_1,\ldots,\vv e_n$ 下的矩阵为 $\vv A$，则 $\mathscr T^*$ 在同一组基下的矩阵为 $\vv A^\herm$.
	\end{theorem}
	\begin{proof}
		固定 $\vv\beta$，$\vv\alpha\mapsto\langle\mathscr T\vv\alpha,\vv\beta\rangle$ 是线性函数. 由引理~\ref{lem:riesz}，有唯一的向量（记为 $\mathscr T^*\vv\beta$）使 $\langle\mathscr T\vv\alpha,\vv\beta\rangle=\langle\vv\alpha,\mathscr T^*\vv\beta\rangle$. 这样定义的 $\mathscr T^*$ 是线性的：对任意 $\vv\alpha$，
		\[\langle\vv\alpha,\mathscr T^*(\vv\beta_1+k\vv\beta_2)\rangle=\langle\mathscr T\vv\alpha,\vv\beta_1\rangle+\bar k\langle\mathscr T\vv\alpha,\vv\beta_2\rangle=\langle\vv\alpha,\mathscr T^*\vv\beta_1+k\mathscr T^*\vv\beta_2\rangle,\]
		由引理中的唯一性即得. 伴随的唯一性也由该唯一性得到.
		
		设 $\mathscr T^*$ 的矩阵为 $\vv B=(b_{ij})$. 由命题~\ref{prop:unitarycoord}，矩阵的 $(i,j)$ 元是第 $j$ 个基向量的像的第 $i$ 个坐标：$a_{ij}=\langle\mathscr T\vv e_j,\vv e_i\rangle$，$b_{ij}=\langle\mathscr T^*\vv e_j,\vv e_i\rangle$. 于是
		\[b_{ij}=\overline{\langle\vv e_i,\mathscr T^*\vv e_j\rangle}=\overline{\langle\mathscr T\vv e_i,\vv e_j\rangle}=\bar a_{ji},\]
		即 $\vv B=\vv A^\herm$.
	\end{proof}
	
	\begin{pitfallbox}
		\begin{enumerate}
			\item ``伴随矩阵''$\vv A^*$（由代数余子式构成，第五章）与``伴随变换''$\mathscr T^*$（对应共轭转置）只是中文名相同，二者毫无关系. 本书用 $\vv A^\herm$ 表示共轭转置，以示区别.
			\item 只有在\textbf{标准正交基}下，伴随变换的矩阵才是共轭转置. 在一般的基下，内积不等于坐标的标准内积，结论不成立.
		\end{enumerate}
	\end{pitfallbox}
	
	\begin{proposition}\label{prop:adjointrules}
		设 $\mathscr S,\mathscr T$ 是酉空间 $V$ 上的线性变换，$k\in\mathbb{C}$. 则
		\begin{enumerate}
			\item $(\mathscr S+\mathscr T)^*=\mathscr S^*+\mathscr T^*$，$(k\mathscr T)^*=\bar k\mathscr T^*$，$(\mathscr{ST})^*=\mathscr T^*\mathscr S^*$，$(\mathscr T^*)^*=\mathscr T$；
			\item $\ker\mathscr T^*=(\imm\mathscr T)^\perp$，$\imm\mathscr T^*=(\ker\mathscr T)^\perp$.
		\end{enumerate}
	\end{proposition}
	\begin{proof}
		(1) 取标准正交基，由定理~\ref{thm:adjoint} 化为共轭转置的运算规则.
		
		(2) $\vv\beta\in\ker\mathscr T^*$ 当且仅当 $\langle\vv\alpha,\mathscr T^*\vv\beta\rangle=0$ 对一切 $\vv\alpha$ 成立，即 $\langle\mathscr T\vv\alpha,\vv\beta\rangle=0$ 对一切 $\vv\alpha$ 成立，即 $\vv\beta\perp\imm\mathscr T$. 把这个结论用于 $\mathscr T^*$，得 $\ker\mathscr T=(\imm\mathscr T^*)^\perp$，两边取正交补即得第二式.
	\end{proof}
	
	对复矩阵，(2) 就是 $\ker\vv A^\herm=(\imm\vv A)^\perp$，$\imm\vv A^\herm=(\ker\vv A)^\perp$，这是第十章四个基本子空间（定理~\ref{thm:fourspaces}）的复数版本：实数情形的转置，在复数情形都要换成共轭转置.
	
	\section{\hC 酉变换与酉矩阵}\label{sec:unitarymatrix}
	
	\begin{definition}
		酉空间 $V$ 上保持内积的线性变换 $\mathscr U$，即对一切 $\vv\alpha,\vv\beta$ 有 $\langle\mathscr U\vv\alpha,\mathscr U\vv\beta\rangle=\langle\vv\alpha,\vv\beta\rangle$，称为\textbf{酉变换}\index{you3 bian4huan4@酉变换}. 满足 $\vv U^\herm\vv U=\vv I$ 的复方阵称为\textbf{酉矩阵}\index{you3 ju3zhen4@酉矩阵}. 若存在酉矩阵 $\vv U$ 使 $\vv B=\vv U^\herm\vv A\vv U$，称 $\vv A$ 与 $\vv B$ \textbf{酉相似}\index{you3 xiang1si4@酉相似}.
	\end{definition}
	
	由推论~\ref{cor:onesided}，$\vv U^\herm\vv U=\vv I$ 意味着 $\vv U^{-1}=\vv U^\herm$，所以酉相似是相似的特例. 实的酉矩阵就是正交矩阵.
	
	\begin{theorem}\label{thm:unitaryeq}
		设 $\mathscr U$ 是酉空间 $V$ 上的线性变换. 以下条件等价：
		\begin{enumerate}
			\item $\mathscr U$ 是酉变换；
			\item $\mathscr U$ 保持长度：$\|\mathscr U\vv\alpha\|=\|\vv\alpha\|$；
			\item $\mathscr U$ 把某一组（从而每一组）标准正交基变为标准正交基；
			\item $\mathscr U$ 在某一组（从而每一组）标准正交基下的矩阵是酉矩阵；
			\item $\mathscr U^*\mathscr U=\mathscr I$，即 $\mathscr U$ 可逆且 $\mathscr U^{-1}=\mathscr U^*$.
		\end{enumerate}
	\end{theorem}
	\begin{proof}
		(1)$\Rightarrow$(2) 取 $\vv\beta=\vv\alpha$. (2)$\Rightarrow$(1)：由极化恒等式 \eqref{eq:polarization}，内积可以用长度表示，保持长度就保持内积. 注意实数情形的二项公式在这里不够用，必须用四项的公式.
		
		(1)$\Leftrightarrow$(3)$\Leftrightarrow$(4) 与定理~\ref{thm:orthequiv} 的证明相同，只需把转置换成共轭转置：$\vv U^\herm\vv U$ 的 $(i,j)$ 元是 $\vv U$ 的第 $j$ 列与第 $i$ 列的标准内积.
		
		(1)$\Leftrightarrow$(5)：$\langle\mathscr U\vv\alpha,\mathscr U\vv\beta\rangle=\langle\mathscr U^*\mathscr U\vv\alpha,\vv\beta\rangle$. 它对一切 $\vv\alpha,\vv\beta$ 等于 $\langle\vv\alpha,\vv\beta\rangle$，当且仅当 $\mathscr U^*\mathscr U\vv\alpha=\vv\alpha$ 对一切 $\vv\alpha$ 成立（引理~\ref{lem:riesz} 中的唯一性）.
	\end{proof}
	
	\begin{proposition}\label{prop:unitaryprop}
		设 $\vv U,\vv U_1,\vv U_2$ 是 $n$ 阶酉矩阵，则
		\begin{enumerate}
			\item $\vv U$ 的各列构成 $\mathbb{C}^n$ 的标准正交基，各行也是；
			\item $\vv U^\herm$，$\vv U^{-1}$，$\vv U_1\vv U_2$ 都是酉矩阵；
			\item $|\det\vv U|=1$；
			\item $\vv U$ 的特征值的模都是 $1$，属于不同特征值的特征向量互相正交.
		\end{enumerate}
	\end{proposition}
	\begin{proof}
		(1)(2) 由 $\vv U^\herm\vv U=\vv U\vv U^\herm=\vv I$ 直接验证. (3) $|\det\vv U|^2=\overline{\det\vv U}\det\vv U=\det(\vv U^\herm\vv U)=1$. (4) 若 $\vv{Ux}=\lambda\vv x$，$\vv x\ne\zero$，则 $\|\vv x\|=\|\vv{Ux}\|=|\lambda|\,\|\vv x\|$，故 $|\lambda|=1$. 再设 $\vv{Uy}=\mu\vv y$，$\mu\ne\lambda$，则
		\[\langle\vv x,\vv y\rangle=\langle\vv{Ux},\vv{Uy}\rangle=\lambda\bar\mu\langle\vv x,\vv y\rangle.\]
		由 $|\mu|=1$ 知 $\bar\mu=\mu^{-1}$，$\lambda\bar\mu=\lambda/\mu\ne1$，所以 $\langle\vv x,\vv y\rangle=0$.
	\end{proof}
	
	\begin{example}
		\begin{enumerate}
			\item 对角矩阵 $\diag(\mathrm e^{\mathrm i\theta_1},\ldots,\mathrm e^{\mathrm i\theta_n})$ 是酉矩阵：它把每个坐标``转''一个角度.
			\item 置换矩阵与正交矩阵都是酉矩阵.
			\item 以 Fourier 基 $\vv f_0,\ldots,\vv f_{n-1}$ 为列的矩阵 $\vv F$ 是酉矩阵（命题~\ref{prop:fourierbasis}）. 离散 Fourier 变换就是 $\vv x\mapsto\vv F^\herm\vv x$，它保持长度——这正是 Parseval 等式.
		\end{enumerate}
	\end{example}
	
	\begin{motivebox}
		第十章（命题~\ref{prop:orthmatrix}）说，正交矩阵的实特征值只能是 $\pm1$；二维旋转的特征值则是一对共轭复数 $\mathrm e^{\pm\mathrm i\theta}$，第十章只能把它们``绕开''. 放到复数域里，正交矩阵是酉矩阵的特例，命题~\ref{prop:unitaryprop} 一次说清：所有特征值都在复平面的单位圆上，$\pm1$ 不过是单位圆与实轴的两个交点. 实数情形看起来零碎的结论，在复数情形是一个整体.
	\end{motivebox}
	
	\section{\hC Hermite 矩阵}\label{sec:hermitian}
	
	\begin{definition}
		满足 $\vv A^\herm=\vv A$ 的复方阵称为 \textbf{Hermite 矩阵}\index{hermite ju3zhen4@Hermite 矩阵}（也称自伴矩阵）；满足 $\vv A^\herm=-\vv A$ 的称为\textbf{反 Hermite 矩阵}. 对 Hermite 矩阵 $\vv A$，函数 $\vv x\mapsto\vv x^\herm\vv A\vv x$（$\vv x\in\mathbb{C}^n$）称为 \textbf{Hermite 二次型}\index{hermite er4ci4xing2@Hermite 二次型}. 若 $\vv x^\herm\vv A\vv x>0$ 对一切 $\vv x\ne\zero$ 成立，称 $\vv A$ 是\textbf{正定} Hermite 矩阵.
	\end{definition}
	
	Hermite 矩阵的对角元满足 $a_{ii}=\bar a_{ii}$，都是实数；关于对角线对称的两个元素互为共轭. 实的 Hermite 矩阵就是实对称矩阵. 由定理~\ref{thm:adjoint}，Hermite 矩阵正是自伴变换在标准正交基下的矩阵.
	
	\begin{theorem}\label{thm:hermitianeig}
		Hermite 矩阵的特征值都是实数，属于不同特征值的特征向量互相正交.
	\end{theorem}
	\begin{proof}
		设 $\vv{Ax}=\lambda\vv x$，$\vv x\ne\zero$. 则 $\lambda\langle\vv x,\vv x\rangle=\langle\vv{Ax},\vv x\rangle=\langle\vv x,\vv{Ax}\rangle=\bar\lambda\langle\vv x,\vv x\rangle$，而 $\langle\vv x,\vv x\rangle>0$，故 $\lambda=\bar\lambda$. 再设 $\vv{Ay}=\mu\vv y$，$\mu\ne\lambda$，则 $\lambda\langle\vv x,\vv y\rangle=\langle\vv{Ax},\vv y\rangle=\langle\vv x,\vv{Ay}\rangle=\bar\mu\langle\vv x,\vv y\rangle=\mu\langle\vv x,\vv y\rangle$（$\mu$ 是实数），故 $\langle\vv x,\vv y\rangle=0$.
	\end{proof}
	
	把这个证明与引理~\ref{lem:symreal} 比较：那里为了证明实对称矩阵的特征值是实数，临时引入了 $\bar{\vv x}^\trans\vv A\vv x$，现在看来，那正是 Hermite 二次型. 复数域才是这个证明的``自然环境''.
	
	\begin{example}\label{eg:hermitian2}
		$\vv A=\begin{pmatrix}2&1-\mathrm i\\1+\mathrm i&3\end{pmatrix}$ 是 Hermite 矩阵. 特征多项式为 $\lambda^2-5\lambda+\big(6-|1-\mathrm i|^2\big)=\lambda^2-5\lambda+4$，特征值为 $1$ 与 $4$，都是实数. 对应的特征向量可取 $\vv x_1=(1-\mathrm i,\,-1)^\trans$ 与 $\vv x_2=(1-\mathrm i,\,2)^\trans$. 验证正交：$\langle\vv x_1,\vv x_2\rangle=(1-\mathrm i)\overline{(1-\mathrm i)}+(-1)\cdot2=2-2=0$.
	\end{example}
	
	\begin{proposition}\label{prop:hermitianform}
		\begin{enumerate}
			\item 复方阵 $\vv A$ 是 Hermite 矩阵，当且仅当 $\vv x^\herm\vv A\vv x$ 对一切 $\vv x\in\mathbb{C}^n$ 都是实数.
			\item 任一复方阵都能唯一地写成 $\vv A=\vv H+\vv K$，其中 $\vv H$ 是 Hermite 矩阵，$\vv K$ 是反 Hermite 矩阵.
		\end{enumerate}
	\end{proposition}
	\begin{proof}
		(1) 若 $\vv A^\herm=\vv A$，则数 $\vv x^\herm\vv A\vv x$ 的共轭（即它的共轭转置）为 $\vv x^\herm\vv A^\herm\vv x=\vv x^\herm\vv A\vv x$，是实数. 反之，设 $\vv x^\herm\vv A\vv x$ 总是实数，则 $\vv x^\herm\vv A\vv x=\vv x^\herm\vv A^\herm\vv x$. 令 $\vv B=\vv A-\vv A^\herm$，则 $\vv x^\herm\vv B\vv x=0$ 对一切 $\vv x$ 成立. 分别取 $\vv x=\vv e_j$，$\vv e_j+\vv e_k$，$\vv e_j+\mathrm i\vv e_k$（$j\ne k$），依次得
		\[b_{jj}=0,\qquad b_{jk}+b_{kj}=0,\qquad\mathrm i(b_{jk}-b_{kj})=0,\]
		所以 $\vv B=\vv O$，$\vv A$ 是 Hermite 矩阵.
		
		(2) 存在性：令 $\vv H=\frac12(\vv A+\vv A^\herm)$，$\vv K=\frac12(\vv A-\vv A^\herm)$. 唯一性：若 $\vv A=\vv H'+\vv K'$ 是另一种写法，则 $\vv H-\vv H'=\vv K'-\vv K$ 既是 Hermite 矩阵又是反 Hermite 矩阵，只能为零.
	\end{proof}
	
	(1) 中用到了 $\vv x=\vv e_j+\mathrm i\vv e_k$，这在实数情形是做不到的. 事实上实数情形的类似结论不成立：$\vv B=\begin{pmatrix}0&1\\-1&0\end{pmatrix}$ 不是零矩阵，但对一切实向量 $\vv x$ 都有 $\vv x^\trans\vv B\vv x=0$. 复数提供了``更多的试验向量''.
	
	(2) 类似于复数的实部与虚部：$z=a+b\mathrm i$ 中，$a=\frac12(z+\bar z)$，$b\mathrm i=\frac12(z-\bar z)$. 反 Hermite 矩阵恰是 $\mathrm i$ 乘以 Hermite 矩阵，于是 $\vv A=\vv H_1+\mathrm i\vv H_2$，$\vv H_1,\vv H_2$ 都是 Hermite 矩阵. Hermite 矩阵扮演了``实数''的角色.
	
	\begin{pitfallbox}
		实 Hermite 矩阵就是实对称矩阵；但\textbf{复对称矩阵}（$\vv A^\trans=\vv A$）一般不是 Hermite 矩阵，第十一章与本节的结论对它都不成立. 例如 $\vv A=\begin{pmatrix}1&\mathrm i\\\mathrm i&-1\end{pmatrix}$ 是复对称矩阵，但 $\vv A^2=\vv O$：它的特征值全为 $0$，而 $\vv A\ne\vv O$，所以它甚至不能对角化. 在复数域中，起作用的是共轭转置，不是转置.
	\end{pitfallbox}
	
	\paragraph*{Rayleigh 商.} 设 Hermite 矩阵 $\vv A$ 的特征值为 $\lambda_1\ge\cdots\ge\lambda_n$. 下一节将证明，$\mathbb{C}^n$ 有一组由 $\vv A$ 的特征向量构成的标准正交基. 承认这一点，与定理~\ref{thm:rayleigh} 完全相同地可得
	\[\max_{\|\vv x\|=1}\vv x^\herm\vv A\vv x=\lambda_1,\qquad\min_{\|\vv x\|=1}\vv x^\herm\vv A\vv x=\lambda_n.\]
	特别地，Hermite 矩阵正定，当且仅当它的特征值都是正数.
	
	\begin{historybox}{从海森堡到冯·诺依曼}
		\sidefig[0.3\linewidth]{image/helgoland.jpg}{黑尔戈兰岛}[fig:helgoland]
		1855 年，法国数学家埃尔米特（C.~Hermite）在研究数论中的二次型时，考察了今天以他命名的``共轭对称''的型，并证明了相应的矩阵的特征值都是实数. 七十年后，这个结论在物理学中找到了归宿. 序言讲过，1925 年海森堡在北海的黑尔戈兰岛上写出了矩阵力学的雏形. 在随后由玻恩、约当和海森堡建立的理论中，位置、动量、能量这些\textbf{可观测量}都用矩阵表示，测量所得的数值是矩阵的特征值. 测量值必须是实数，所以这些矩阵必须是 Hermite 矩阵——定理~\ref{thm:hermitianeig} 成了物理理论的一条前提. 不同测量结果对应的状态互相正交，也是这个定理的一部分. 1932 年，冯·诺依曼（J.~von~Neumann）在《量子力学的数学基础》中，把量子力学建立在无穷维的复内积空间（Hilbert 空间）及其上的自伴算子之上，这一框架沿用至今.
	\end{historybox}
	
	\section{\hC 复矩阵的分解定理}\label{sec:complexdecomp}
	
	\subsection{Schur 分解}
	
	第十一章的谱定理只对实对称矩阵成立. 对任意复方阵，我们退一步，不要求化为对角矩阵，只要求化为上三角矩阵.
	
	\begin{theorem}[Schur 分解]\label{thm:schur}\index{schur fen1jie3@Schur 分解}
		每个复方阵 $\vv A$ 都酉相似于一个上三角矩阵：存在酉矩阵 $\vv U$，使 $\vv U^\herm\vv A\vv U=\vv T$ 是上三角矩阵. $\vv T$ 的对角元是 $\vv A$ 的全部特征值.
	\end{theorem}
	\begin{proof}
		对阶数 $n$ 归纳，$n=1$ 显然. 由代数基本定理，$\vv A$ 的特征多项式有复根 $\lambda_1$，取对应的单位特征向量 $\vv u_1$. 由 Gram–Schmidt 正交化，把 $\vv u_1$ 扩充为 $\mathbb{C}^n$ 的标准正交基，以它们为列作酉矩阵 $\vv U_1=(\vv u_1,\vv U')$. 则
		\[\vv U_1^\herm\vv A\vv U_1=\begin{pmatrix}\vv u_1^\herm\vv A\vv u_1&\vv u_1^\herm\vv A\vv U'\\\vv U'^\herm\vv A\vv u_1&\vv U'^\herm\vv A\vv U'\end{pmatrix}=\begin{pmatrix}\lambda_1&*\\\zero&\vv B\end{pmatrix},\]
		其中左下块 $\vv U'^\herm(\lambda_1\vv u_1)=\zero$，因为 $\vv U'$ 的列与 $\vv u_1$ 正交. 由归纳假设，有 $n-1$ 阶酉矩阵 $\vv U_2$ 使 $\vv U_2^\herm\vv B\vv U_2$ 为上三角矩阵. 令 $\vv U=\vv U_1\diag(1,\vv U_2)$，它是酉矩阵之积，$\vv U^\herm\vv A\vv U$ 为上三角矩阵. 相似矩阵有相同的特征多项式，上三角矩阵的特征值就是它的对角元.
	\end{proof}
	
	这个证明与谱定理（定理~\ref{thm:spectral}）的证明几乎逐字相同. 区别在两处：第一，复方阵总有特征值，而实矩阵不一定有实特征值；第二，对称矩阵的右上块自动为零，一般矩阵则不然，所以只能得到上三角矩阵.
	
	\begin{motivebox}
		为什么实数域上没有这样一般的结论？归纳的第一步需要一个特征向量. 实矩阵 $\begin{pmatrix}0&-1\\1&0\end{pmatrix}$（旋转 $90^\circ$）没有实特征值，第一步就走不下去. 实数域上能得到的最好结果是``分块上三角''，对角线上是一阶或二阶的块——又是``实数世界里，不可约的东西至多是二维的''（第七章、第十章）. 复数域保证了特征值的存在，Schur 分解才能对一切方阵成立.
	\end{motivebox}
	
	\begin{example}
		第九章的 $\vv A=\begin{pmatrix}3&1\\-1&1\end{pmatrix}$ 不能对角化，Jordan 标准形为 $\vv J_2(2)$. 求它的 Schur 分解：特征值 $2$ 对应的单位特征向量为 $\vv u_1=\frac1{\sqrt2}(1,-1)^\trans$，扩充为标准正交基，取 $\vv u_2=\frac1{\sqrt2}(1,1)^\trans$. 令 $\vv U=(\vv u_1,\vv u_2)$，计算得
		\[\vv U^\herm\vv A\vv U=\begin{pmatrix}2&2\\0&2\end{pmatrix}.\]
		与 Jordan 标准形相比，Schur 分解的上三角部分没有那么``干净''（右上角是 $2$ 而不是 $1$），但所用的 $\vv U$ 是酉矩阵. 酉矩阵保持长度，计算误差不会被放大；Jordan 标准形所需的过渡矩阵则可能非常``病态''（第九章的计算实践）. 这就是数值计算中用 Schur 分解代替 Jordan 标准形的原因.
	\end{example}
	
	\subsection{正规矩阵的谱定理}
	
	什么样的复方阵能被酉矩阵对角化？若 $\vv A=\vv U\vv D\vv U^\herm$，$\vv D$ 是对角矩阵，则 $\vv A^\herm=\vv U\bar{\vv D}\vv U^\herm$，而对角矩阵 $\vv D$ 与 $\bar{\vv D}$ 可交换，所以
	\[\vv A\vv A^\herm=\vv U\vv D\bar{\vv D}\vv U^\herm=\vv U\bar{\vv D}\vv D\vv U^\herm=\vv A^\herm\vv A.\]
	这就引出了下面的定义.
	
	\begin{definition}
		满足 $\vv A\vv A^\herm=\vv A^\herm\vv A$ 的复方阵称为\textbf{正规矩阵}\index{zheng4gui1 ju3zhen4@正规矩阵}(normal matrix).
	\end{definition}
	
	Hermite 矩阵、反 Hermite 矩阵、酉矩阵都是正规矩阵：对它们，$\vv A^\herm$ 分别等于 $\vv A$、$-\vv A$、$\vv A^{-1}$，都与 $\vv A$ 可交换.
	
	\begin{theorem}[正规矩阵的谱定理]\label{thm:normalspectral}\index{zheng4gui1 ju3zhen4 de pu3 ding4li3@正规矩阵的谱定理}
		复方阵 $\vv A$ 能被酉矩阵对角化，当且仅当 $\vv A$ 是正规矩阵. 换言之，$\mathbb{C}^n$ 有一组由 $\vv A$ 的特征向量构成的标准正交基，当且仅当 $\vv A\vv A^\herm=\vv A^\herm\vv A$.
	\end{theorem}
	\begin{proof}
		必要性上面已经证明. 充分性：由 Schur 分解，$\vv A=\vv U\vv T\vv U^\herm$，$\vv T$ 是上三角矩阵. 由 $\vv A$ 正规，
		\[\vv T\vv T^\herm=\vv U^\herm\vv A\vv A^\herm\vv U=\vv U^\herm\vv A^\herm\vv A\vv U=\vv T^\herm\vv T,\]
		即 $\vv T$ 也是正规矩阵. 下面证明正规的上三角矩阵是对角矩阵. 比较 $\vv T\vv T^\herm$ 与 $\vv T^\herm\vv T$ 的 $(1,1)$ 元：前者是 $\vv T$ 的第一行的长度平方 $\sum_{j}|t_{1j}|^2$，后者是 $\vv T$ 的第一列的长度平方 $|t_{11}|^2$（第一列除 $t_{11}$ 外都是零）. 二者相等，所以 $t_{12}=\cdots=t_{1n}=0$. 于是 $\vv T=\diag(t_{11},\vv T_1)$，$\vv T_1$ 仍是正规的上三角矩阵，对它重复同样的论证，即知 $\vv T$ 是对角矩阵.
	\end{proof}
	
	\begin{corollary}\label{cor:normalclasses}
		Hermite 矩阵、酉矩阵、反 Hermite 矩阵都能被酉矩阵对角化. 一个正规矩阵是 Hermite 矩阵、酉矩阵、反 Hermite 矩阵，当且仅当它的特征值分别全是实数、全在单位圆上、全是纯虚数（包括 $0$）.
	\end{corollary}
	\begin{proof}
		设 $\vv A=\vv U\vv D\vv U^\herm$，$\vv D=\diag(\lambda_1,\ldots,\lambda_n)$. 则 $\vv A^\herm=\vv U\bar{\vv D}\vv U^\herm$，$\vv A^\herm=\vv A$ 当且仅当 $\bar{\vv D}=\vv D$，即 $\lambda_i$ 都是实数. 其余两种情形同理：$\vv A^\herm\vv A=\vv I$ 当且仅当 $|\lambda_i|^2=1$；$\vv A^\herm=-\vv A$ 当且仅当 $\bar\lambda_i=-\lambda_i$.
	\end{proof}
	
	\begin{example}
		实矩阵 $\vv A=\begin{pmatrix}1&-1\\1&1\end{pmatrix}$ 不是对称矩阵，第十一章的谱定理不适用，它也确实不能用正交矩阵对角化（它没有实特征值）. 但 $\vv A\vv A^\trans=\vv A^\trans\vv A=2\vv I$，它是正规矩阵. 特征值为 $1\pm\mathrm i$，对应的单位特征向量为
		\[\vv u_1=\frac1{\sqrt2}(1,-\mathrm i)^\trans,\qquad\vv u_2=\frac1{\sqrt2}(1,\mathrm i)^\trans,\]
		$\langle\vv u_1,\vv u_2\rangle=\frac12\big(1+(-\mathrm i)\cdot\overline{\mathrm i}\big)=\frac12(1+\mathrm i^2)=0$. 令 $\vv U=(\vv u_1,\vv u_2)$，得 $\vv U^\herm\vv A\vv U=\diag(1+\mathrm i,1-\mathrm i)$. 从几何上看，$\vv A=\sqrt2\begin{pmatrix}\cos45^\circ&-\sin45^\circ\\\sin45^\circ&\cos45^\circ\end{pmatrix}$ 是``旋转 $45^\circ$ 再放大 $\sqrt2$ 倍''，特征值 $1+\mathrm i=\sqrt2\,\mathrm e^{\mathrm i\pi/4}$ 恰好记录了旋转的角度与放大的倍数.
	\end{example}
	
	由第九章的观点看，正规矩阵的 Jordan 标准形中每个 Jordan 块都是一阶的. 但``能对角化''与``能被酉矩阵对角化''是两回事：前者要求有 $n$ 个线性无关的特征向量，后者要求它们还能取成两两正交的.
	
	\begin{codebox}{Schur 分解}
		\mstep{目标} 用库函数求 Schur 分解，观察正规矩阵与非正规矩阵的区别.
		
		\mstep{准备} \texttt{numpy}；\texttt{scipy.linalg.schur}（取参数 \texttt{output='complex'}）.
		
		\mstep{任务}
		\begin{enumerate}
			\item 对 $\vv A=\begin{pmatrix}3&1\\-1&1\end{pmatrix}$ 求 Schur 分解 $\vv A=\vv U\vv T\vv U^\herm$，验证 $\vv U^\herm\vv U=\vv I$ 与 $\vv T$ 是上三角矩阵，与正文的例子对照.
			\item 随机生成一个 $5\times5$ 复矩阵 $\vv B$，令 $\vv H=\vv B+\vv B^\herm$. 分别求 $\vv B$ 与 $\vv H$ 的 Schur 分解，打印 $\vv T$ 的严格上三角部分的 Frobenius 范数.
			\item 对 $\vv B$ 计算 $\|\vv B\vv B^\herm-\vv B^\herm\vv B\|$，作为``偏离正规''的程度，与任务 2 中的范数比较.
		\end{enumerate}
		
		\mstep{检查点} 任务 1 中 $|t_{12}|=2$（$\vv U$ 的列可以相差一个模为 $1$ 的因子，$t_{12}$ 的辐角可能不同）；任务 2 中 $\vv H$ 的严格上三角部分的范数在 $10^{-15}$ 量级，$\vv B$ 的则明显不为零.
		
		\mstep{思考} Schur 分解中的 $\vv U$ 不唯一（特征值的排列顺序可以不同）. $\vv T$ 的严格上三角部分的 Frobenius 范数是否与 $\vv U$ 的选取有关？（提示：$\|\vv T\|_F=\|\vv A\|_F$，而对角部分是特征值.）
	\end{codebox}
	
	\subsection{复矩阵的奇异值分解}
	
	第十一章的奇异值分解（定理~\ref{thm:svd}）的证明，只用到 $\vv A^\trans\vv A$ 是半正定的实对称矩阵，以及谱定理. 对复矩阵，$\vv A^\herm\vv A$ 是 Hermite 矩阵，且 $\vv x^\herm\vv A^\herm\vv A\vv x=\|\vv{Ax}\|^2\ge0$，由推论~\ref{cor:normalclasses} 它能被酉矩阵对角化，特征值非负. 把原证明中的转置都换成共轭转置，就得到：
	
	\begin{theorem}[复矩阵的奇异值分解]\label{thm:complexsvd}
		设 $\vv A$ 是 $m\times n$ 复矩阵，$\rk\vv A=r$. 则存在 $m$ 阶酉矩阵 $\vv U$、$n$ 阶酉矩阵 $\vv V$，以及 $\sigma_1\ge\cdots\ge\sigma_r>0$，使 $\vv A=\vv U\vv\Sigma\vv V^\herm$，其中 $\vv\Sigma$ 的形状与定理~\ref{thm:svd} 相同. $\sigma_i$ 是 $\vv A^\herm\vv A$ 的非零特征值的算术平方根.
	\end{theorem}
	
	\paragraph*{*极分解.} 设 $\vv A$ 是 $n$ 阶复方阵，$\vv A=\vv U\vv\Sigma\vv V^\herm$. 在中间插入 $\vv V^\herm\vv V=\vv I$，
	\[\vv A=(\vv U\vv V^\herm)(\vv V\vv\Sigma\vv V^\herm)=\vv W\vv P,\]
	其中 $\vv W=\vv U\vv V^\herm$ 是酉矩阵，$\vv P=\vv V\vv\Sigma\vv V^\herm$ 是特征值非负的 Hermite 矩阵（半正定）. 这称为 $\vv A$ 的\textbf{极分解}\index{ji2 fen1jie3@极分解}，它类比于复数的极坐标表示 $z=\mathrm e^{\mathrm i\theta}r$：酉矩阵相当于``模为 $1$ 的复数''，半正定 Hermite 矩阵相当于``非负实数''. 几何上，每个线性变换都是先沿某组正交方向伸缩（$\vv P$），再作一个保持长度的变换（$\vv W$）.
	
	\subsection{实数与复数的对照}
	
	本章的主要结论与第十、十一章的对应关系如下.
	\begin{table}[htbp]
		\centering
		\begin{tabular}{lll}
			\toprule
			& 实数域 & 复数域 \\
			\midrule
			标准内积 & $\vv y^\trans\vv x$ & $\vv y^\herm\vv x$ \\
			保持内积的矩阵 & 正交矩阵 $\vv Q^\trans\vv Q=\vv I$ & 酉矩阵 $\vv U^\herm\vv U=\vv I$ \\
			自伴的矩阵 & 实对称矩阵 & Hermite 矩阵 \\
			能被``好的基''对角化 & 实对称矩阵（正交对角化） & 正规矩阵（酉对角化） \\
			任意方阵能化成 & ——（不一定有实特征值） & 上三角矩阵（Schur 分解） \\
			奇异值分解 & $\vv U\vv\Sigma\vv V^\trans$ & $\vv U\vv\Sigma\vv V^\herm$ \\
			\bottomrule
		\end{tabular}
		\caption{实内积空间与酉空间的对照}
		\label{tab:realcomplex}
	\end{table}
	
	表中``实数域''一列的每一条，都是``复数域''一列的特例. 反过来，从实数域推广到复数域，所需的改动几乎只有一条：把转置换成共轭转置.
	
	\begin{modelbox}{用离散 Fourier 变换做季节调整}
		\mstep{问题} 某地连续 $48$ 个月的月度用电量有明显的季节性：夏、冬两季高，春、秋两季低. 如何去掉季节性，看出用电量本身的变化？
		
		\mstep{假设} 数据可以拆成三部分：季节成分（以 $12$ 个月为周期，一年内平均为零）、长期成分与随机波动.
		
		\mstep{模型} 把 $48$ 个月的数据记为 $\vv x=(x_0,\ldots,x_{47})^\trans\in\mathbb{C}^{48}$（数据是实数，看作复向量），在 Fourier 基 $\vv f_0,\ldots,\vv f_{47}$ 下展开，$\vv x=\sum c_k\vv f_k$. 设季节成分 $\vv s$ 满足 $s_{m+12}=s_m$. 把求和按年分成四段，它的 Fourier 坐标为
		\[\langle\vv s,\vv f_k\rangle=\frac1{\sqrt{48}}\Big(\sum_{m=0}^{11}s_m\bar\omega^{km}\Big)\big(1+\bar\omega^{12k}+\bar\omega^{24k}+\bar\omega^{36k}\big),\qquad\omega=\mathrm e^{2\pi\mathrm i/48}.\]
		$\bar\omega^{12k}=\mathrm e^{-\pi\mathrm ik/2}$，后一个括号是公比为 $\mathrm e^{-\pi\mathrm ik/2}$ 的等比数列之和：当 $4\mid k$ 时为 $4$，否则为 $0$. 再由``一年内平均为零''，$k=0$ 的坐标也为零. 所以\textbf{季节成分只出现在 $k=4,8,\ldots,44$ 这 $11$ 个频率上}.
		
		\mstep{求解} 计算 $\vv x$ 的离散 Fourier 变换，把 $k=4,8,\ldots,44$ 的坐标置为零，其余不变，再变换回去（$\vv x\mapsto\sum c_k\vv f_k$），得到的向量就是季节调整后的序列. 这一步就是把 $\vv x$ 正交投影到 $\spn\{\vv f_k\mid 4\nmid k\ \text{或}\ k=0\}$ 上.
		
		\mstep{解读}
		\begin{itemize}
			\item 季节调整是一次正交投影：去掉的部分与保留的部分正交，由勾股定理，$\|\vv x\|^2$ 恰好分成``季节部分''与``其余部分''两块，可以用二者之比衡量季节性的强弱.
			\item 调整后的序列适合作环比（与上个月比）；未调整的序列只适合作同比（与去年同月比）——同比本身就消去了以 $12$ 个月为周期的成分.
		\end{itemize}
		
		\mstep{局限}
		\begin{itemize}
			\item 长期增长趋势不是周期的，它在 Fourier 基下的坐标会``泄漏''到所有频率上，包括 $k=4,8,\ldots$. 实际做法通常先估计并去掉趋势，再处理季节性.
			\item 春节的日期在公历中逐年移动，它的影响不严格以 $12$ 个月为周期. 统计部门使用的方法（如 X-13）还要专门处理节假日、交易日和异常值.
		\end{itemize}
	\end{modelbox}
	
	\newpage
